\documentclass[11pt]{article}
\usepackage[margin=1in]{geometry}
\usepackage[T1]{fontenc}
\usepackage{graphicx}
\usepackage{hyperref}
\usepackage{enumitem}
\setlength{\parskip}{0.5em}
\setlength{\parindent}{0pt}

\title{state space models for sequence modeling}
\author{arXiv GraphRAG --- phase-6 deep research}
\date{\today}

\begin{document}
\maketitle

\begin{abstract}
An excellent, comprehensive review that masterfully connects historical RNN stability issues to the modern selective SSM frontier with high technical precision.
\end{abstract}

\section{From Recurrence to Selection: A Survey of State Space Models in Sequence Modeling}

Sequence modeling serves as the cornerstone of modern artificial intelligence, enabling machines to process temporal and ordered data across domains such as natural language processing, signal analysis, and genomics. Historically, the field has been defined by a fundamental tension between computational efficiency and modeling capacity. While early architectures provided the groundwork for understanding temporal dependencies, the subsequent rise of the Transformer architecture introduced a significant efficiency gap, where the quadratic scaling of attention mechanisms limits the ability to process increasingly long sequences. Resolving this bottleneck is critical for the development of next-generation models that require both high-fidelity long-range modeling and the linear-time inference necessary for massive-scale deployment.

This review traces the trajectory of sequence modeling from its foundations in recurrent neural networks and early Transformers through the subsequent attempts to bridge the efficiency gap using sparse and linear attention mechanisms. We then examine the emergence of structured state space models as a transformative paradigm, detailing their evolution into selective and recurrent architectures that refine how information is gated and maintained over time. By exploring the theoretical unification and mathematical formalism that underpin these systems, this survey provides a cohesive framework for understanding their underlying mechanics. Finally, we look toward the frontiers of the field, discussing the development of hardware-aware, multi-scale, and specialized SSMs that aim to redefine the computational limits of sequence processing.

\subsection{Foundations of Sequence Modeling: RNNs and Early Transformers}

Early sequence modeling was dominated by Recurrent Neural Networks (RNNs), which faced significant stability challenges during training. A primary obstacle was the tendency for gradients to either vanish or explode; researchers addressed this by introducing gradient norm clipping and soft constraints to manage gradient stability (\texttt{arXiv:1211.5063}). To specifically combat these issues, Unitary Evolution RNNs utilized a unitary weight matrix with eigenvalues of absolute value exactly 1, employing complex-domain hidden states to facilitate the learning of very long-term dependencies (\texttt{arXiv:1511.06464}). Further refinements in architecture led to the development of Gated Recurrent Neural Networks, such as LSTMs and GRUs, which utilized gating mechanisms to manage sequence handling (\texttt{arXiv:1412.3555}).

While gated architectures improved performance in domains like speech and music modeling (\texttt{arXiv:1412.3555}), they remained computationally constrained by their sequential nature. Efforts to increase parallelism led to the development of Quasi-Recurrent Neural Networks, which replaced traditional recurrence with alternating convolutional layers and a minimalist recurrent pooling function (\texttt{arXiv:1611.01576}). This approach achieved up to 16 times faster training and testing speeds compared to LSTMs while maintaining predictive accuracy across language modeling and sentiment classification tasks (\texttt{arXiv:1611.01576}).

A fundamental paradigm shift occurred with the introduction of the Transformer, which dispensed with recurrence and convolutions entirely in favor of multi-head self-attention mechanisms (\texttt{arXiv:1706.03762}). This architecture's novelty lies in its ability to draw global dependencies through attention alone, making it significantly more parallelizable and requiring less training time than its recurrent or convolutional predecessors (\texttt{arXiv:1706.03762}). This shift enabled the rise of large-scale pre-trained models. BERT introduced a deep bidirectional Transformer encoder that jointly conditions on both left and right context in all layers (\texttt{arXiv:1810.04805}). Subsequent research focused on optimizing these foundations: RoBERTa demonstrated that BERT was significantly undertrained and that optimized hyperparameters and larger datasets could exceed existing performance (\texttt{arXiv:1907.11692}), while T5 established a unified text-to-text framework to convert all language problems into a consistent format (\texttt{arXiv:1910.10683}). Throughout these advancements, the Adam optimizer—utilizing adaptive estimates of first and second-order gradient moments—remained a critical tool for the efficient optimization of these large-scale parameters (\texttt{arXiv:1412.6980}).

\subsection{The Efficiency Gap: Sparse and Linear Attention Mechanisms}

To address the quadratic computational complexity of standard Transformers, several architectural innovations have emerged, primarily focusing on sparse attention patterns or linear-time approximations. One foundational approach involves segment-level recurrence; Transformer-XL (\texttt{arXiv:1901.02860}) mitigates context fragmentation by enabling dependencies 80\% longer than RNNs and 450\% longer than vanilla Transformers through a novel positional encoding scheme and recurrence mechanism. While this improves context handling, it does not fundamentally alter the underlying attention complexity.

In contrast, sparse attention mechanisms aim to reduce complexity by restricting the attention matrix. Sparse Transformers (\texttt{arXiv:1904.10509}) utilize sparse factorizations to achieve \$O(n \textbackslash{}sqrt\{n\})\$ complexity, while Routing Transformers (\texttt{arXiv:2003.05997}) employ an online k-means routing module to achieve \$O(n\textasciicircum{}\{1.5\}d)\$ complexity, allowing for training on sequences up to 8192 tokens. Further refinements in sparsity include the Longformer (\texttt{arXiv:2004.05150}), which scales linearly by combining local windowed attention with task-motivated global attention, and Big Bird (\texttt{arXiv:2007.14062}), which achieves linear scaling by incorporating \$O(1)\$ global tokens. Big Bird specifically distinguishes itself by remaining a universal approximator and Turing complete. Reformer (\texttt{arXiv:2001.04451}) offers a different path to efficiency, utilizing locality-sensitive hashing (LSH) to achieve \$O(L \textbackslash{}log L)\$ time complexity, paired with reversible residual layers to minimize the memory footprint.

A third category of research focuses on linearizing the attention mechanism entirely. Linformer (\texttt{arXiv:2006.04768}) reduces complexity to \$O(n)\$ via low-rank matrix approximation. Performer (\texttt{arXiv:2009.14794}) approaches this through the FAVOR+ mechanism, which uses positive Orthogonal Random features to provide an unbiased estimation of the softmax attention kernel, achieving linear time and space complexity. This concept is extended by the "Transformers are RNNs" framework (\texttt{arXiv:2006.16236}), which expresses self-attention as a linear dot-product of kernel feature maps. This allows for an iterative implementation via the associativity property of matrix products, enabling \$O(N)\$ complexity and accelerating autoregressive prediction by up to 4000x. However, the novelty of this linearization via kernel feature maps has faced critique regarding its specific implementation of \$O(N)\$ relations (\texttt{arXiv:2406.09062}).

These diverse efficient architectures are often evaluated using the Long Range Arena (LRA) benchmark (\texttt{arXiv:2011.04006}), which provides a unified framework for testing models on sequences ranging from 1K to 16K tokens across text, image, and mathematical domains. Together, these developments represent a shift from the rigid quadratic scaling of original Transformers toward more flexible, scalable, and recurrent-like structures.

\subsection{The Emergence of Structured State Space Models}

The evolution of state space models (SSMs) for sequence modeling is characterized by a transition from theoretical frameworks for memory to highly efficient, structured architectures capable of handling extremely long-range dependencies. This progression began with the HiPPO framework (\texttt{arXiv:2008.07669}), which introduced a method for online function approximation by projecting continuous signals onto polynomial bases. This provided a mechanism for recurrent neural networks to compress cumulative history with timescale robustness, establishing a foundation for subsequent structured models (\texttt{arXiv:2008.07669}).

Building upon this, the Structured State Space (S4) model (\texttt{arXiv:2111.00396}) introduced a novel parameterization of the state transition matrix \$A\$ using a Normal Plus Low-Rank (NPLR) structure. This innovation allowed for stable diagonalization and efficient computation via Cauchy kernel reduction, enabling the model to solve the 16k-length Path-X task with 88\% accuracy—a feat previously impossible for other models (\texttt{arXiv:2111.00396}). While S4 achieved 60× faster generation than Transformers on WikiText-103, it remains susceptible to numerical instability if \$A\$ possesses eigenvalues in the right half-plane (\texttt{arXiv:2111.00396}). Furthermore, ablation studies suggest that the performance gains may be driven more by HiPPO initialization than the NPLR parameterization itself (\texttt{arXiv:2111.00396}).

Subsequent research focused on simplifying these structures. The Linear State-Space Layer (LSSL) generalized RNNs and temporal convolutions through a linear continuous-time representation (\texttt{arXiv:2110.13985}). This trajectory toward simplification culminated in the realization that diagonal state matrices could match S4's performance without complex low-rank corrections (\texttt{arXiv:2203.14343}). Specifically, the S4D model (\texttt{arXiv:2206.11893}) demonstrated that restricting the state matrix to a diagonal form—while systematically initializing it to recover the S4 kernel in the infinite limit—maintains effectiveness across image, audio, and medical time-series domains (\texttt{arXiv:2206.11893}).

Efficiency and expressivity were further addressed through architectural refinements. The Simplified State Space Layers (\texttt{arXiv:2208.04933}) replaced multiple single-input, single-output (SISO) models with a single multi-input, multi-output (MIMO) layer, utilizing parallel scans to achieve 98.5\% accuracy on Path-X (\texttt{arXiv:2208.04933}). To bridge the expressivity gap with attention mechanisms in language modeling, the H3 model (\texttt{arXiv:2212.14052}) combined SSMs with a novel state-passing algorithm and FlashConv, scaling nearly linearly with sequence length (\texttt{arXiv:2212.14052}). Finally, the theoretical capacity of these models was reinforced by proofs that SSMs can encode n-gram rules by restricting the spectrum of the state transition matrix (\texttt{arXiv:2306.17184}).

\subsection{Evolution of Selective and Recurrent Architectures}

The evolution of state space models (SSMs) is characterized by a tension between the parallelizable training efficiency of Transformers and the constant-complexity inference of Recurrent Neural Networks (RNNs). RWKV addresses this by combining linear attention with RNN-style inference, allowing for parallelizable training while maintaining constant computational and memory complexity during inference for NLP tasks \texttt{arXiv:2305.13048}. While RWKV reconciles these two paradigms, subsequent research has focused on enhancing the selectivity and stability of state transitions.

A significant shift occurred with the introduction of Mamba, which moves away from Linear Time Invariant (LTI) constraints by utilizing selective SSMs. In this architecture, parameters (\$\textbackslash{}Delta, B,\$ and \$C\$) are functions of the input, enabling content-based reasoning through selective propagation or forgetting of information \texttt{arXiv:2312.00752}. This selectivity allows Mamba to handle sequences up to \$10\textasciicircum{}6\$ tokens and achieve \$5\textbackslash{}times\$ higher generation throughput than Transformers \texttt{arXiv:2312.00752}. However, this departure from LTI-based convolutional efficiency introduces a trade-off: it necessitates complex, hardware-specific implementations, such as fused kernels and recomputation, to prevent \$O(BLD)\$ complexity issues \texttt{arXiv:2312.00752}.

Other research trajectories have focused on the mathematical stability and structural formulation of these transitions. To combat the "curse of memory"—specifically the exponential decay of memory in standard SSMs—StableSSM introduces stable reparameterization techniques to enhance optimization stability \texttt{arXiv:2311.19694}. This concept of stability is further refined in S7, which utilizes stable reparameterization to enable input-dependent filtering while simultaneously controlling gradient norms for long-sequence modeling \texttt{arXiv:2410.03464}. While Mamba achieves selectivity through input-dependent parameters, S7 explicitly integrates this selectivity with stability-focused reparameterization to manage the dynamics of state transitions \texttt{arXiv:2410.03464}.

Alternative approaches have explored different mathematical foundations for state transitions. ConvSSM utilizes parallel scans applied to convolutional recurrences to achieve subquadratic parallelization, establishing an equivalence between ConvSSM dynamics and SSMs for spatiotemporal modeling \texttt{arXiv:2310.19694}. In contrast, Spectral SSMs utilize fixed convolutional filters via a spectral filtering algorithm, providing provable robustness for long-range dependencies without the need for learning the filters themselves \texttt{arXiv:2312.06837}. Together, these developments represent a move from rigid, linear architectures toward highly adaptive, hardware-aware, and mathematically stabilized systems capable of processing massive sequences across diverse modalities, from genomics to dynamical systems \texttt{arXiv:2312.00752}, \texttt{arXiv:2403.16899}.

\subsection{Theoretical Unification and Mathematical Formalism}

The theoretical landscape of state space models (SSMs) has evolved from structured recurrence toward complex architectures designed to achieve linear scaling while managing long-range dependencies (\texttt{arXiv:2503.18970}). A primary advancement in this unification is the Structured State Space Duality (SSD) framework, which establishes a formal mathematical connection between SSMs and attention variants through the use of structured semiseparable matrices, providing the theoretical foundation for the Mamba-2 architecture (\texttt{arXiv:2405.21060}).

Mathematical formalisms for memory management further differentiate current research trajectories. One approach focuses on selective gating mechanisms that dynamically filter input relevance, utilizing information-theoretic tools—specifically mutual information and rate-distortion theory—to achieve efficient memory compression (\texttt{arXiv:2410.03158}). In contrast, other frameworks conceptualize SSM design through the lens of online learning. In these models, state transition rules are not merely gated but are derived from closed-form solutions to online associative recall problems, facilitating linear decoding efficiency (\texttt{arXiv:2407.14207}).

A critical theoretical tension exists regarding the necessity of time-varying dynamics for achieving selectivity. While the progression from S4 to Mamba emphasizes integrated structured recurrence to mitigate vanishing gradients and quadratic memory complexity (\texttt{arXiv:2503.18970}), recent work challenges the assumption that selectivity requires breaking Linear Time-Invariant (LTI) properties. The Geometric SSM utilizes geometric control theory and a dynamic residual generator to excite distinct invariant subspaces, enabling selective modeling while maintaining the efficiency of LTI systems (\texttt{arXiv:2505.17932}). This approach offers a significant strength in its ability to recognize multi-token patterns, achieving over 99\% accuracy on extended induction head tasks where Mamba’s memoryless selection mechanism fails (\texttt{arXiv:2505.17932}). However, a notable weakness is that this geometric approach has only been validated on synthetic benchmarks and sMNIST, lacking the large-scale language modeling or complex natural language dataset validation required to prove its general utility (\texttt{arXiv:2505.17932}).

To provide a rigorous basis for comparing these diverse mathematical approaches, recent research has introduced synthetic benchmarking frameworks. These frameworks utilize controllable memory functions to evaluate how different architectures capture distinct temporal structures, providing a standardized method for investigating sequence modeling theory (\texttt{arXiv:2506.05678}).

\subsection{Frontiers: Hardware, Multi-Scale, and Specialized SSMs}

Recent advancements in state-space models (SSMs) have shifted toward specialized hardware implementations, multi-scale representations, and complex temporal dynamics to overcome the limitations of standard linear recurrences. On the hardware front, research bifurcates between specialized accelerators and memory-centric architectures. For instance, the EpochCore accelerator utilizes a systolic array with LIMA-PE and ProDF to optimize structured SSMs for long-range tasks (\texttt{arXiv:2507.21394}), whereas the compute-in-memory (CIM) approach leverages memristors to enable fully asynchronous, energy-efficient processing for event-driven sequences (\texttt{arXiv:2511.13912}).

To address the limitations of single-scale modeling, MS-SSM introduces a multi-resolution framework using an input-dependent scale-mixer (\texttt{arXiv:2512.23824}). This architecture demonstrates significant performance gains in hierarchical tasks, such as achieving 63.04\% on ListOps—nearly doubling Mamba's 38.02\% (\texttt{arXiv:2512.23824}). However, while MS-SSM provides empirical strength, it lacks formal theoretical justification for its scale-dependent initialization and provides only qualitative complexity claims rather than detailed FLOPs or parameter comparisons (\texttt{arXiv:2512.23824}). In contrast, Prefix-Scannable Models (PSMs) provide a more rigorous generalization of SSMs by allowing arbitrary state aggregation via parallel prefix scans (\texttt{arXiv:2506.10918}). PSMs achieve \$O(1)\$ amortized compute per token and \$\textbackslash{}log(N)\$ memory for sequence length \$N\$, offering a different path to efficiency through operator flexibility rather than resolution-based fusion (\texttt{arXiv:2506.10918}).

Specialized dynamics also represent a growing frontier for modeling non-stationary environments. While standard SSMs often rely on input-dependent gating, time-varying SSMs utilize a dictionary of basis functions to parameterize \$A[t]\$, \$B[t]\$, and \$C[t]\$ matrices (\texttt{arXiv:2605.15311}). This allows for modeling switching dynamics in non-stationary signals, such as achieving 19.6 dB SI-SNR in speech denoising compared to 14.2 dB for time-invariant models (\texttt{arXiv:2605.15311}). A notable weakness of this approach is that stability guarantees rely on the simplifying assumption that \$A[t]\$ is diagonal (\texttt{arXiv:2605.15311}). This pursuit of dynamic expressivity is mirrored in the coefficient dynamics framework, which unifies SSMs, RNNs, and Transformers by treating output coefficients as outputs of autonomous linear dynamical systems (\texttt{arXiv:2510.09389}).

Finally, SSM principles are being extended to non-sequential domains and adaptive learning. The MP-SSM framework integrates SSM computation into message-passing neural networks, ensuring permutation equivariance and efficient long-range propagation for graph learning (\texttt{arXiv:2505.18728}). Similarly, the L2T-Hyena architecture enhances the Hyena SSM through a Learning-to-Teach paradigm, using a teacher model to dynamically adapt the student's loss function by trading off the primary cross-entropy objective against a regularization term (\texttt{arXiv:2511.05926}).

\subsection{Tensions and Open Problems}

A primary tension in the literature concerns the fundamental architectural trade-off between the global expressive power of attention mechanisms and the computational efficiency of recurrent or convolutional structures. On one hand, the Transformer paradigm, as established in \texttt{arXiv:1706.03762}, provides a robust framework for modeling complex dependencies through all-to-all attention. However, this approach faces significant scaling challenges regarding sequence length, a limitation that \texttt{arXiv:1901.02860} attempts to mitigate by introducing segment-level recurrence. Recent surveys such as \texttt{arXiv:2412.11211} argue that the entire Transformer lineage—including large-scale few-shot learners like \texttt{arXiv:2005.14165} and efficient variants like \texttt{arXiv:2001.04451}—is fundamentally constrained by the quadratic complexity of the attention mechanism, suggesting that the field must move toward more scalable alternatives to handle truly long-range dependencies.

A second debate centers on the mathematical formulation of State Space Models (SSMs) and the necessity of structural complexity. Early work on HiPPO in \texttt{arXiv:2008.07669} emphasized the importance of optimal polynomial projections to maintain long-term memory. This view has been challenged by the finding in \texttt{arXiv:2203.14343} that diagonal state spaces can achieve comparable effectiveness to more complex structured models, suggesting a potential path toward simplification. This tension extends to the implementation of convolutions; while \texttt{arXiv:2102.02611} proposes continuous kernels for sequential data, more recent critiques in \texttt{arXiv:2302.06646} argue that simpler, hardware-efficient long convolutions offer superior practical performance for sequence modeling.

Finally, there is an emerging disagreement regarding the "selectivity" of state-space architectures versus their inherent stability. The introduction of Mamba in \texttt{arXiv:2312.00752} marked a shift toward selective state spaces that can adapt to input content, a move intended to overcome the limitations of Linear Time-Invariant (LTI) models. However, recent critical analyses in \texttt{arXiv:2503.18970} suggest that the evolution from S4 to Mamba involves complex trade-offs that may compromise the theoretical guarantees of earlier SSMs. This has led to a fragmented landscape where researchers must choose between the highly adaptive, content-aware mechanisms of selective models and the more predictable, mathematically grounded dynamics of traditional state-space frameworks.

\subsection{Outlook}

The frontier of sequence modeling is rapidly shifting from the pursuit of mere scaling toward the integration of multi-scale dynamics and hardware-aware design. We are seeing a transition from monolithic architectures to hybrid systems, such as the multi-scale approaches in \texttt{arXiv:2512.23824} and the adaptive "learn-to-teach" frameworks in \texttt{arXiv:2511.05926}. As the field matures, the focus is moving beyond algorithmic complexity toward the physical realities of deployment, evidenced by recent work on systolic array accelerators (\texttt{arXiv:2507.21394}) and compute-in-memory implementations (\texttt{arXiv:2511.13912}). The next generation of models will likely be defined by their ability to unify sequential-parallel duality (\texttt{arXiv:2506.05678}) with geometric and graph-based inductive biases, creating architectures that are as efficient in hardware as they are expressive in theory.

\section*{References}
\begin{enumerate}\small
  \item On the difficulty of training Recurrent Neural Networks (2012). \texttt{arXiv:1211.5063}
  \item Adam: A Method for Stochastic Optimization (2014). \texttt{arXiv:1412.6980}
  \item Learning Phrase Representations using RNN Encoder-Decoder for
  Statistical Machine Translation (2014). \texttt{arXiv:1406.1078}
  \item Neural Machine Translation by Jointly Learning to Align and Translate (2014). \texttt{arXiv:1409.0473}
  \item Empirical Evaluation of Gated Recurrent Neural Networks on Sequence
  Modeling (2014). \texttt{arXiv:1412.3555}
  \item Deep Residual Learning for Image Recognition (2015). \texttt{arXiv:1512.03385}
  \item Unitary Evolution Recurrent Neural Networks (2015). \texttt{arXiv:1511.06464}
  \item Layer Normalization (2016). \texttt{arXiv:1607.06450}
  \item Gaussian Error Linear Units (GELUs) (2016). \texttt{arXiv:1606.08415}
  \item Pointer Sentinel Mixture Models (2016). \texttt{arXiv:1609.07843}
  \item Language Modeling with Gated Convolutional Networks (2016). \texttt{arXiv:1612.08083}
  \item Quasi-Recurrent Neural Networks (2016). \texttt{arXiv:1611.01576}
  \item The LAMBADA dataset: Word prediction requiring a broad discourse context (2016). \texttt{arXiv:1606.06031}
  \item Attention Is All You Need (2017). \texttt{arXiv:1706.03762}
  \item Decoupled Weight Decay Regularization (2017). \texttt{arXiv:1711.05101}
  \item BERT: Pre-training of Deep Bidirectional Transformers for Language
  Understanding (2018). \texttt{arXiv:1810.04805}
  \item GLUE: A Multi-Task Benchmark and Analysis Platform for Natural Language
  Understanding (2018). \texttt{arXiv:1804.07461}
  \item An Empirical Evaluation of Generic Convolutional and Recurrent Networks
  for Sequence Modeling (2018). \texttt{arXiv:1803.01271}
  \item Think you have Solved Question Answering? Try ARC, the AI2 Reasoning
  Challenge (2018). \texttt{arXiv:1803.05457}
  \item Generating Long Sequences with Sparse Transformers (2019). \texttt{arXiv:1904.10509}
  \item Transformer-XL: Attentive Language Models Beyond a Fixed-Length Context (2019). \texttt{arXiv:1901.02860}
  \item RoBERTa: A Robustly Optimized BERT Pretraining Approach (2019). \texttt{arXiv:1907.11692}
  \item Exploring the Limits of Transfer Learning with a Unified Text-to-Text
  Transformer (2019). \texttt{arXiv:1910.10683}
  \item Compressive Transformers for Long-Range Sequence Modelling (2019). \texttt{arXiv:1911.05507}
  \item Transformers are RNNs: Fast Autoregressive Transformers with Linear
  Attention (2020). \texttt{arXiv:2006.16236}
  \item HiPPO: Recurrent Memory with Optimal Polynomial Projections (2020). \texttt{arXiv:2008.07669}
  \item Language Models are Few-Shot Learners (2020). \texttt{arXiv:2005.14165}
  \item Reformer: The Efficient Transformer (2020). \texttt{arXiv:2001.04451}
  \item Long Range Arena: A Benchmark for Efficient Transformers (2020). \texttt{arXiv:2011.04006}
  \item Longformer: The Long-Document Transformer (2020). \texttt{arXiv:2004.05150}
  \item Rethinking Attention with Performers (2020). \texttt{arXiv:2009.14794}
  \item Linformer: Self-Attention with Linear Complexity (2020). \texttt{arXiv:2006.04768}
  \item An Image is Worth 16x16 Words: Transformers for Image Recognition at
  Scale (2020). \texttt{arXiv:2010.11929}
  \item Big Bird: Transformers for Longer Sequences (2020). \texttt{arXiv:2007.14062}
  \item Efficient Content-Based Sparse Attention with Routing Transformers (2020). \texttt{arXiv:2003.05997}
  \item Efficient Transformers: A Survey (2020). \texttt{arXiv:2009.06732}
  \item Efficiently Modeling Long Sequences with Structured State Spaces (2021). \texttt{arXiv:2111.00396}
  \item Combining Recurrent, Convolutional, and Continuous-time Models with
  Linear State-Space Layers (2021). \texttt{arXiv:2110.13985}
  \item CKConv: Continuous Kernel Convolution For Sequential Data (2021). \texttt{arXiv:2102.02611}
  \item On the Parameterization and Initialization of Diagonal State Space
  Models (2022). \texttt{arXiv:2206.11893}
  \item FlashAttention: Fast and Memory-Efficient Exact Attention with
  IO-Awareness (2022). \texttt{arXiv:2205.14135}
  \item Diagonal State Spaces are as Effective as Structured State Spaces (2022). \texttt{arXiv:2203.14343}
  \item Hungry Hungry Hippos: Towards Language Modeling with State Space Models (2022). \texttt{arXiv:2212.14052}
  \item Simplified State Space Layers for Sequence Modeling (2022). \texttt{arXiv:2208.04933}
  \item It's Raw! Audio Generation with State-Space Models (2022). \texttt{arXiv:2202.09729}
  \item How to Train Your HiPPO: State Space Models with Generalized Orthogonal
  Basis Projections (2022). \texttt{arXiv:2206.12037}
  \item Convolutional State Space Models for Long-Range Spatiotemporal Modeling (2023). \texttt{arXiv:2310.19694}
  \item Spectral State Space Models (2023). \texttt{arXiv:2312.06837}
  \item Block-State Transformers (2023). \texttt{arXiv:2306.09539}
  \item Effectively Modeling Time Series with Simple Discrete State Spaces (2023). \texttt{arXiv:2303.09489}
  \item Simple Hardware-Efficient Long Convolutions for Sequence Modeling (2023). \texttt{arXiv:2302.06646}
  \item Mamba: Linear-Time Sequence Modeling with Selective State Spaces (2023). \texttt{arXiv:2312.00752}
  \item Structured State Space Models for In-Context Reinforcement Learning (2023). \texttt{arXiv:2303.03982}
  \item Resurrecting Recurrent Neural Networks for Long Sequences (2023). \texttt{arXiv:2303.06349}
  \item StableSSM: Alleviating the Curse of Memory in State-space Models through
  Stable Reparameterization (2023). \texttt{arXiv:2311.14495}
  \item State space models can express n-gram languages (2023). \texttt{arXiv:2306.17184}
  \item RWKV: Reinventing RNNs for the Transformer Era (2023). \texttt{arXiv:2305.13048}
  \item State-Space Modeling in Long Sequence Processing: A Survey on Recurrence in the Transformer Era (2024). \texttt{arXiv:2406.09062}
  \item Learning Long Sequences in Spiking Neural Networks (2024). \texttt{arXiv:2401.00955}
  \item Mathematical Formalism for Memory Compression in Selective State Space
  Models (2024). \texttt{arXiv:2410.03158}
  \item State Space Models as Foundation Models: A Control Theoretic Overview (2024). \texttt{arXiv:2403.16899}
  \item Deep Learning-based Approaches for State Space Models: A Selective
  Review (2024). \texttt{arXiv:2412.11211}
  \item S7: Selective and Simplified State Space Layers for Sequence Modeling (2024). \texttt{arXiv:2410.03464}
  \item Mamba-360: Survey of State Space Models as Transformer Alternative for
  Long Sequence Modelling: Methods, Applications, and Challenges (2024). \texttt{arXiv:2404.16112}
  \item Slot State Space Models (2024). \texttt{arXiv:2406.12272}
  \item Longhorn: State Space Models are Amortized Online Learners (2024). \texttt{arXiv:2407.14207}
  \item LongVQ: Long Sequence Modeling with Vector Quantization on Structured
  Memory (2024). \texttt{arXiv:2404.11163}
  \item Transformers are SSMs: Generalized Models and Efficient Algorithms
  Through Structured State Space Duality (2024). \texttt{arXiv:2405.21060}
  \item MS-SSM: A Multi-Scale State Space Model for Efficient Sequence Modeling (2025). \texttt{arXiv:2512.23824}
  \item Advancing Intelligent Sequence Modeling: Evolution, Trade-offs, and Applications of State- Space Architectures from S4 to Mamba (2025). \texttt{arXiv:2503.18970}
  \item Geometric SSM: LTI State Space Models for Selective Tasks (2025). \texttt{arXiv:2505.17932}
  \item Message-Passing State-Space Models: Improving Graph Learning with Modern Sequence Modeling (2025). \texttt{arXiv:2505.18728}
  \item Design Principles for Sequence Models via Coefficient Dynamics (2025). \texttt{arXiv:2510.09389}
  \item L2T-Hyena: Enhancing State-Space Models with an Adaptive Learn-to-Teach Framework (2025). \texttt{arXiv:2511.05926}
  \item Compute-in-Memory Implementation of State Space Models for Event Sequence Processing (2025). \texttt{arXiv:2511.13912}
  \item Sequential-Parallel Duality in Prefix Scannable Models (2025). \texttt{arXiv:2506.10918}
  \item SeRpEnt: Selective Resampling for Expressive State Space Models (2025). \texttt{arXiv:2501.11729}
  \item Systolic Array-based Accelerator for Structured State-Space Models (2025). \texttt{arXiv:2507.21394}
  \item Numerical Investigation of Sequence Modeling Theory using Controllable Memory Functions (2025). \texttt{arXiv:2506.05678}
  \item Time-Varying Deep State Space Models for Sequences with Switching Dynamics (2026). \texttt{arXiv:2605.15311}
\end{enumerate}

\end{document}
